% ---
% id: protocol-compiler/thm:measurements-and-statistics
% kind: thm
% title: "Exact measurement families and the statistics functional"
% summary: >-
%   Declared measurement families are exact POVMs block by block, per factor
%   and, in POVM form, per Alice sector in the general format; the product and
%   denoted ensemble families are POVMs by theorem. So D(iii) holds, and the
%   statistics are a distribution on every normalised positive block-list state.
% prominence: standard
% status: accepted
% uses: [def:compile, def:well-formed-spec, def:exact, def:phi, thm:typing, lem:hermitian-constraint-values, def:feasible, def:weighted, def:validator, def:declaration, def:target, thm:sifting-and-key]
% mentions: [thm:chain-consumption, lem:weighted, thm:validator-soundness, foundations/req:exact-emission, req:derived-objects]
% assumes: [A-SPEC]
% refs: []
% provenance: newly-derived
% implements: certiqs_key.protocol_compiler.validation:check_v3
% lean:
% notes: "The \\eqref{eq:predcount} of the proof of (iii) is to an equation of thm:validator-soundness, a mention: the sentence counts decisions and proves nothing. A-SPEC is declared because the proof applies thm:typing and thm:sifting-and-key, whose hypothesis it is; A-ALICE-MEAS is not, no step reading the measurement Alice performs."
% ---

\begin{theorem}\label{thm:measurements-and-statistics}
Let $\Spec$ be a \tref{well-formed specification}
(Definition~\ref{def:well-formed-spec}) satisfying
Assumption~\ref{ass:spec}, and
let $\Decl$ be emitted by Definition~\ref{def:compile} in the
\tref{basic format}. Then:
\begin{enumerate}[label=(\roman*),leftmargin=2.2em]
\item $\{M_x\}_{x\in\Xset}$ is a POVM on $A$ and, for each $s\in\SecSet$, the
block family $\{F_v^{(s)}\}_{v\in\Vset}$ is a POVM on $\Qsys_s$ --- each
predicate holding \emph{of the declared numbers}, decided exactly --- and each
of the completeness sums $\sum_xM_x$ and $\sum_vF_v^{(s)}$, formed in the
declared order, rounds nowhere, so that each computed sum \emph{is} the sum it
is written as (Definition~\ref{def:exact});
\item every $F_v$ and every $\hat O_k$ is sector-block-diagonal
\emph{by construction}: it is emitted as the list
$\big(F_v^{(s)}\big)_{s}$, respectively $\big(\hat O_k^{(s)}\big)_{s}$, and
the operator it denotes is $\bigoplus_sF_v^{(s)}$. The two formulations of
block-diagonality remain equivalent --- for an operator $F$ on $\Qsys$,
$\Pi_sF\Pi_{s'}=0$ for all $s\ne s'$ if and only if
$F=\bigoplus_s\Pi_sF\Pi_s$ --- and that equivalence is what licenses the
\emph{type}: a list of blocks denotes exactly the operators the vanishing
condition picks out, and no other. It is not checked, because the condition
has no data to be checked on;
\item hence D(iii) holds: its typing half by (ii) and check V1; its first
four predicates --- those on $\{M_x\}$ and on each block family
$\{F_v^{(s)}\}_v$ --- by (i); and its fifth, Hermiticity of
$\hat O_k^{(s)}$ for each $(k,s)$, by the last clause of check V3
(Definition~\ref{def:validator}), which decides it exactly on the emitted
constraint entries and which (i) does not state; and
\item $P_{\Decl}[\omega]$ of \eqref{eq:jointdist} and $\Phimap[\omega]$ are
affine in $\omega$, $\Phimap[\omega]$ depends on $\omega$ only through
$P_{\Decl}[\omega]$, and the fibres of $\phi\circ a$ partition
$\Xset\times\Vset$, so $\Phimap[\omega]$ is the pushforward of
$P_{\Decl}[\omega]$ and inherits its total mass. For every block list $\omega$
with $\omega^{(s)}\succeq0$ for each $s$ and $\tr\omega=1$ --- in particular
for every $\omega\in\Feas$ --- $P_{\Decl}[\omega]$ is a probability
distribution on $\Xset\times\Vset$ and $\Phimap[\omega]$ one on $\Chatset$.
\end{enumerate}

\emph{The general format.} Let $\Decl$ instead be emitted by the
general-format clauses of Definition~\ref{def:compile}, in the general format
of Definition~\ref{def:declaration}\,(DX1)--(DX6). Then (i)--(iv) hold with
the following readings, of which the four above are the case of the basic
format.
\begin{enumerate}[label=(\roman*$'$),leftmargin=2.6em]
\item for every \tref{measured factor} $f$ and every $s_f\in\SecSet^{f}$ the
block family $\{F^{f,(s_f)}_{v_f}\}_{v_f\in\Vset^{f}}$ is a POVM on
$\Qsys_{f,s_f}$; and for every Alice sector $t$, in \tref{POVM form}
$\{M^{(t)}_x\}_{x\in\Xset}$ is a POVM on $A_t$, and in \tref{ensemble form}
every $N^{(t)}_x$ is Hermitian and positive semidefinite, with
$\sum_xN^{(t)}_x=\tau^{(t)}_A$ and $\tau^{(t)}_A\ne0$ --- each predicate
holding of the declared numbers, decided exactly, and each of these sums,
formed in its declared order, rounding nowhere. The product family
\eqref{eq:productfamily} is a POVM on $\Qsys$ and, in ensemble form, the
family \eqref{eq:ensemblepovm} is a POVM on $A$, both as identities of real
numbers, and neither is formed, decided or summed;
\item every $F^{f}_{v_f}$ is a block list against factor $f$'s sector map and
every $\hat O_k$ one against the sector map of its scope, and $\tau_A$, every
operator of Alice's measurement slot and every general $\Theta_k$ is a block
list along $\mathbb{T}$, each block diagonal by construction as in (ii); and
the product family is block diagonal for the product sector map,
$\Pi_sF_v\Pi_{s'}=0$ for $s\ne s'$, with the $s$-block of
\eqref{eq:productfamily};
\item hence D(iii) holds in its general-format reading
(Definition~\ref{def:target}); and
\item (iv) holds for every list $\omega$ over the
\tref[joint sector]{joint sectors} with every block positive semidefinite and
$\tr\omega=1$ --- in particular for every $\omega\in\Feas$ --- with $F_v$ the
product \eqref{eq:productfamily} and, in ensemble form, $M_x$ the operator
\eqref{eq:ensemblepovm} of the declaration $\Decl$ denotes: $P_{\Decl}[\omega]$
is a probability distribution on $\Xset\times\Vset$ and $\Phimap[\omega]$ one
on $\Chatset$.
\end{enumerate}
\end{theorem}

\begin{proof}
(i) The predicates are decided on the emitted objects, not inferred
(\xref{foundations/req:exact-emission}): check V3 of
Definition~\ref{def:validator} decides, exactly, $M_x=M_x^{\dagger}$ and
$M_x\succeq0$ for each $x$, $\sum_xM_x=\idon{A}$ with the sum rounding
nowhere, and, per block, $F_v^{(s)}=F_v^{(s)\dagger}$, $F_v^{(s)}\succeq0$ and
$\sum_vF_v^{(s)}=\idon{\Qsys_s}$ with that sum rounding nowhere; a declaration
that is emitted at all is one for which V3 did not refuse. Two constructions in
this contract would have licensed an inference instead, and neither is used as
one. Where the specification presents a family as choice-weighted,
Lemma~\ref{lem:weighted}\,(i) shows the weighted family is a POVM whenever the
conditional families are; the emitted family's predicates are nonetheless
decided, because the emitted matrices are the products $q_tN_{u|t}$ as
represented, which the compiler forms exactly and then decides to be binary64
data before emitting them at all (Definition~\ref{def:weighted}), and because
the lemma says nothing about condition~(S) on the emitted family's own sum.
Where $\{F_v\}$ arrives as the target POVM of the last link of $\pcChain$,
completeness is likewise decided by V3 on the emitted blocks and never inferred
from the physical family; Theorem~\ref{thm:chain-consumption} records what the chain is and
is not asked to guarantee. No floating-point budget is therefore needed
anywhere below: nothing here depends on how much rounding a construction
introduced, because the result is decided.

(ii) By Definition~\ref{def:compile} the emitted $\{F_v\}$ and $\{\hat O_k\}$
are block lists against $\Smap$, so the first sentence is the emission rule.
For the equivalence: any $F\in\Lin{\Qsys}$ satisfies
$F=\sum_{s,s'}\Pi_sF\Pi_{s'}$, since $\sum_s\Pi_s=\idon{\Qsys}$ and the $\Pi_s$
are mutually orthogonal projectors; hence $\Pi_sF\Pi_{s'}=0$ for all $s\ne s'$
holds if and only if $F=\sum_s\Pi_sF\Pi_s=\bigoplus_s\Pi_sF\Pi_s$, i.e.\ if and
only if $F$ is the direct sum of its diagonal blocks. So the block lists denote
exactly the block-diagonal operators, which is what makes the type faithful:
nothing expressible as a list violates the condition, and nothing satisfying
the condition fails to be expressible as a list. This is a statement about the
representation, not a check on data. Where the specification presents
$\{F_v\}$ as a choice-weighted family, the weighting is applied blockwise,
$N_{u,t}^{(s)}=q_tN_{u|t}^{(s)}$, so the emitted object is a block list
whatever the presentation (Lemma~\ref{lem:weighted}\,(ii)).

(iii) is (i) and (ii) together with D(iii)'s statement. Its typing half is
(ii), and the sector-map agreement that half rests on is check V1
(Theorem~\ref{thm:typing}, D(i)). Of D(iii)'s five predicates, the four on
$\{M_x\}$ and on the block families $\{F_v^{(s)}\}_v$ are (i). The fifth is
not among them and is cited from where it is decided: D(iii) requires in
addition that $\hat O_k^{(s)}$ be Hermitian for each $(k,s)$, which (i) does
not state, and which the last clause of check V3
(Definition~\ref{def:validator}) decides exactly on the emitted constraint
entries --- a declaration that is emitted at all being one for which V3 did
not refuse, which is the route (i)'s own four take
(Theorem~\ref{thm:typing}). It is among the decisions \eqref{eq:predcount}
counts, as that equation's $m\,|\SecSet|$ term, and it is the hypothesis of
Lemma~\ref{lem:hermitian-constraint-values}, which makes
\eqref{eq:feas-modes} well posed.

(iv) Affinity of both maps is immediate from linearity of the trace in
$\omega$; that $\Phimap[\omega]$ depends on $\omega$ only through
$P_{\Decl}[\omega]$ is the definition \eqref{eq:jointdist}. Since $\phi\circ a$
is a total map $\Xset\times\Vset\to\Chatset$ (Theorem~\ref{thm:typing}, D(i)), the
fibres $\{(x,v):\phi(a(x,v))=\hat c\}$ partition $\Xset\times\Vset$, so
$\Phimap[\omega]$ is the pushforward of $P_{\Decl}[\omega]$ and has the same
total mass.

For the distribution statement, let $\omega$ be a block list with
$\omega^{(s)}\succeq0$ for each $s$ and $\tr\omega=1$. By (i) each $M_x$ and
each $F_v^{(s)}$ is Hermitian with non-negative spectrum, hence
$M_x\succeq0$ and $F_v=\bigoplus_sF_v^{(s)}\succeq0$, positivity of a
block-diagonal operator being positivity of each of its blocks; and
$\omega=\bigoplus_s\omega^{(s)}\succeq0$ for the same reason. Then
non-negativity and realness together, since $M_x\otimes F_v\succeq0$ and
$\omega\succeq0$ give $\tr[(M_x\otimes F_v)\omega]\ge0$ --- the trace of a
product of two positive semidefinite operators is a non-negative \emph{real}
number. And normalization:
\begin{equation*}
  \sum_{x,v}P_{\Decl}[\omega](x,v)
  =\tr\!\Big[\Big(\sum_xM_x\otimes\sum_vF_v\Big)\omega\Big]
  =\tr[(\idon{A}\otimes\idon{\Qsys})\,\omega]=\tr\omega=1 ,
\end{equation*}
by bilinearity of the tensor product and the two completeness identities of
(i). This step is where condition~(S) is used: what (i) delivers is
$\sum_xM_x=\idon{A}$ and $\sum_vF_v^{(s)}=\idon{\Qsys_s}$ as identities of the
declared numbers, and it delivers them only because the sums forming them round
nowhere --- a computed sum equal to the identity would not license this
display. Hence $P_{\Decl}[\omega]$ is a probability distribution, and so is its
pushforward $\Phimap[\omega]$.

\emph{The general format.} (i$'$) Check V3 in its general-format reading
(Definition~\ref{def:validator}) decides, exactly: per measured factor,
$F^{f,(s_f)}_{v_f}$ Hermitian and positive semidefinite for every
$(v_f,s_f)$ and, for each $s_f$,
$\sum_{v_f}F^{f,(s_f)}_{v_f}=\idon{\Qsys_{f,s_f}}$ with its declared-order sum
rounding nowhere; on Alice's data, block by block along $\mathbb{T}$, in POVM
form $M^{(t)}_x$ Hermitian and positive semidefinite for every $(x,t)$ and
$\sum_xM^{(t)}_x=\idon{A_t}$ for every $t$ with that sum rounding nowhere,
and in ensemble form $N^{(t)}_x$ Hermitian and positive semidefinite for
every $(x,t)$ and $\sum_xN^{(t)}_x=\tau^{(t)}_A$ entry by entry with that sum
rounding nowhere; and check V2 in its general-format reading decides
$\tau^{(t)}_A\ne0$ for every $t$ where the slot is in ensemble form. A
declaration that is emitted at all is one for which neither check refused, so
each of these holds of the declared numbers. The product family: every
$F^{f}_{v_f}=\bigoplus_{s_f}F^{f,(s_f)}_{v_f}$ is then Hermitian and positive
semidefinite and $\sum_{v_f}F^{f}_{v_f}=\idon{\Qsys_f}$, an identity of the
declared numbers because each of its blocks' sums rounds nowhere; a tensor
product of Hermitian positive semidefinite operators is Hermitian and
positive semidefinite, and
$\sum_{v}F_v=\sum_{v}\bigotimes_fF^{f}_{v_f}=\bigotimes_f\sum_{v_f}F^{f}_{v_f}
=\bigotimes_f\idon{\Qsys_f}=\idon{\Qsys}$ by multilinearity, the sum running
over $\Vset=\prod_f\Vset^{f}$ (Definition~\ref{def:declaration}\,(DX3)). The
denoted family: Theorem~\ref{thm:sifting-and-key}\,(v), whose hypothesis is the
conditions Definition~\ref{def:exact} states for the ensemble form, which are
those just decided. Neither family is a datum
(Requirement~\ref{req:derived-objects}): nothing of either is formed, decided
or summed, so no summation condition arises for either, and both statements
are identities of the real numbers the declared data denote.

(ii$'$) By the general-format clauses of Definition~\ref{def:compile} each
$F^{f}_{v_f}$ is emitted as a block list against factor $f$'s map and each
$\hat O_k$ against the map of its scope, and by
Definition~\ref{def:declaration}\,(DX2) Alice's data are block lists along
$\mathbb{T}$; for each such list the argument of (ii), with the projectors
$\Pi^{f}_{s_f}$, respectively $\Pi^{A}_t$, in place of $\Pi_s$, gives the
equivalence. For the product family, $\Pi_s=\bigotimes_f\Pi^{f}_{s_f}$ by
\eqref{eq:productmap}, so
$\Pi_sF_v\Pi_{s'}=\bigotimes_f\Pi^{f}_{s_f}F^{f}_{v_f}\Pi^{f}_{s'_f}$; for
$s\ne s'$ some factor has $s_f\ne s'_f$, and its tensor factor vanishes because
$F^{f}_{v_f}$ is block diagonal for factor $f$'s map; and
$\Pi_sF_v\Pi_s=\bigotimes_fF^{f,(s_f)}_{v_f}$ is the block of
\eqref{eq:productfamily}.

(iii$'$) D(iii) in its general-format reading (Definition~\ref{def:target})
asks for the typing and for the predicates of Definition~\ref{def:exact} in
their general-format reading, per measured factor and per Alice sector, with
the ensemble-form conditions in place of those on $\{M_x\}$ where that form is
declared, and it refuses ensemble form in untrusted-source mode. The typing
is (ii$'$), the agreement of each list's map with the map of its factor being
decided by check V1 in its general-format reading (Theorem~\ref{thm:typing},
D(i)); the predicates and summation conditions on the per-factor families and
on Alice's data are (i$'$); Hermiticity of $\hat O^{(s_{f_k})}_k$ for every
entry and every block of its scope is the last clause of check V3 in its
general-format reading; and an ensemble-form slot beside a marginal slot
holding $\nomarg$ is refused by that check's clause on which slot is filled.
A declaration that is emitted at all is one for which V1 and V3 did not
refuse.

(iv$'$) Affinity and the pushforward are as in (iv), with $a$ the function
\eqref{eq:transcript} derives, total into $\Kann$
(Theorem~\ref{thm:sifting-and-key}\,(iv)), and $\phi$ total on $\Kann$
(Theorem~\ref{thm:typing}, D(i)). For the distribution statement let $\omega$
be such a list. $M_x\succeq0$: in POVM form $M_x=\bigoplus_tM^{(t)}_x$ with
every block positive semidefinite by (i$'$), and in ensemble form by
Theorem~\ref{thm:sifting-and-key}\,(v); $F_v\succeq0$ by (i$'$); and
$\omega=\bigoplus_{(t,s)}\omega^{(t,s)}\succeq0$, positivity of a
block-diagonal operator being positivity of each block. So
$\tr[(M_x\otimes F_v)\omega]$ is a non-negative real number. And
$\sum_xM_x=\bigoplus_t\sum_xM^{(t)}_x=\bigoplus_t\idon{A_t}=\idon{A}$, in POVM
form an identity of the declared numbers by (i$'$) and in ensemble form an
identity of real numbers by Theorem~\ref{thm:sifting-and-key}\,(v), while
$\sum_vF_v=\idon{\Qsys}$ by (i$'$); so the normalization display of (iv) holds
with these operators and gives $\sum_{x,v}P_{\Decl}[\omega](x,v)=\tr\omega=1$.
Condition~(S) is used there per measured factor and, in POVM form, per Alice
sector, which is what (i$'$) delivers; no sum over the product alphabet and no
sum of the denoted family is computed, so none needs a summation condition.
\end{proof}
